\section{Methodology}
\label{sec:methodology}

Light transport in layered materials, especially with an arbitrary number of layers, is complex to model efficiently. In this section we will present the statistical approach of \textcite{belcour2018} and our implementation of it and its extensions for anisotropic interfaces \parencite{yamaguchi2019,weier2020}. This approach is capable of estimating light interaction accurately and efficiently, without the need for parameter or material dependent precomputation. The model can be used in offline as well as real-time renderers. \textcite{belcour2018} introduced a statistical framework with atomic operators and a new adding-doubling algorithm to evaluate multiple scattering.

\subsection{Statistics of an Interface}
\label{sec:interface_statistics}

The statistical framework of \textcite{belcour2018} is build upon the property that there exists a bijective mapping between the BSDF parameters of the GGX microfacet model \parencite{walter2007} and its directional statistics. This allows us to estimate the energy $e$, mean $\mu$ and variance $\sigma$ of a BSDF lobe for each layer.\footnote{We will omit the square on $\sigma$ for better readability.} These statistics can then be used while evaluation to instantiate the individual lobes of the resulting BSDF.

%

The GGX BRDF $\rho_{GGX}(\omega_i, \omega_o)$ is defined by the following equations \parencite{walter2007,heitz2014}:

\begin{equation} \label{eq:ggx_brfd}
        \rho_{GGX}(\omega_i, \omega_o) = \frac{F(\dotproduct{\omega_i}{h}) G(\omega_i, \omega_o) D(h)}{4 \dotproduct{\omega_i}{n} \dotproduct{\omega_o}{n}}\\
\end{equation}

where $F$ is the Fresnel term, G is the masking-shadowing term and D is the normal distribution function. For the anisotropic case the equations are parametrized by the two roughness parameters $\alpha_x$ and $\alpha_y$:
\begin{equation} \label{eq:g_d_definition}
    \begin{aligned}
        D(\omega) &= \frac{1}{\pi \alpha_x \alpha_y \left( \frac{\omega_x^2}{\alpha_x^2} + \frac{\omega_y^2}{\alpha_y^2} + \omega_z^2 \right)^2} \\
        G(\omega_i, \omega_o) &= \frac{1}{1 + \Lambda(\omega_i) + \Lambda(\omega_o)} \\
        \text{with }\\
        \Lambda(\omega) &= \frac{-1 + \sqrt{1 + \frac{\alpha_x^2 \omega_x^2 + \alpha_y^2 \omega_y^2}{\omega_z^2}}}{2}
    \end{aligned}
\end{equation}
The Fresnel equations describe the amount of light reflected from a surface. For unpolarized light at the boundary between a dielectric and a conductor the Fresnel reflectance is given by \parencite{pharr2023}:
\begin{equation}\label{eq:fresnel_conductor}
    \begin{aligned}
        F &= \frac{1}{2} \left(r_\parallel^2 + r_\perp^2\right) \\
        \text{with}\\ 
        r_\perp &= \frac{a^2 + b^2 - 2a \cos{\theta} + \cos^2{\theta}}{a^2 + b^2 + 2a \cos{\theta} + \cos^2{\theta}}\\
        r_\parallel &= r_\perp \frac{\cos^2{\theta} (a^2+b^2) - 2a \cos{\theta} \sin^2{\theta} + \sin^4{\theta}}{\cos^2{\theta} (a^2+b^2) + 2a \cos{\theta} \sin^2{\theta} + \sin^4{\theta}}\\
        \text{where}\\
        a^2 + b^2 &= \sqrt{(\eta^2 - \kappa^2 - \sin^2{\theta})^2 + 4\eta^2\kappa^2}
    \end{aligned}
\end{equation}
$\eta + i \kappa$ is the ratio of the complex indices of refraction of the two media and $\theta$ is the angle between the incident direction and the surface normal. At the interface of two dielectric media the equations can be simplified.

%

When studying the outgoing distribution of directions of a GGX lobe orthographically projected to the $(x,y)$ plane it can be noticed that this results in approximately a circular shape (see Fig. \ref{fig:projective_space}). Using this assumption the mean can be parametrized with respect to $\begin{pmatrix} u & v \end{pmatrix} = \begin{pmatrix} \omega_x & \omega_y \end{pmatrix}$ of directions $\omega = \begin{pmatrix} \omega_x & \omega_y & \omega_z \end{pmatrix}$ and variance with just a scalar. This introduces some inaccuracy as GGX lobes are noticeably skewed at grazing angles, however as GGX lobes are only parametrized by an energy, incident direction and roughness, higher order statistics cannot be incorporated anyway.
\begin{figure}[!htb]
    \includegraphics[width = \columnwidth, keepaspectratio]{fromBelcour/projectivespace.png}
    \caption{The statistical analysis is performed in a projective space as GGX lobes projected to the base plane can be approximated using a circular distribution. Image source: \textcite{belcour2018}}
    \label{fig:projective_space}
\end{figure}

%

In order to track the variance $\sigma$ correctly, a mapping from GGX roughness $\alpha$ to a linear variance space is required. Therefore, a invertable transformation function $\sigma = f(\alpha)$ is needed. A reflection at a rough surface increases the variance not linearly with respect to roughness. To accurately account for multiple bounces the following property should hold: $\sigma_{12} = f(\alpha_1) + f(\alpha_2)$ for two bounces at surfaces with roughness $\alpha_1$ and $\alpha_2$. Numerically evaluating this, \textcite{belcour2018} found the mapping:
\begin{equation} \label{eq:variance_mapping}
    \begin{aligned}
        f(\alpha) &= \frac{\alpha^{1.1}}{\left(1-\alpha^{1.1}\right)} \\
        f^{-1}(\sigma) &= \left(\frac{\sigma}{1+\sigma}\right)^{\frac{1}{1.1}}
    \end{aligned}
\end{equation}
This mapping does not fit the linear space perfectly but is close for roughnesses below $0.4$. \textcite{weier2020} improved this mapping by using an optimizer to find the parameters of a log sigmoid:
\begin{equation} \label{eq:variance_mapping_improved}
    \begin{aligned}
        f(\alpha) &= \log\left(1+\frac{b \alpha^a}{1- \alpha^a}\right)\\
        f^{-1}(\sigma) &= \left(\frac{e^\sigma - 1}{e^\sigma - 1 + b}\right)^{\frac{1}{a}}
    \end{aligned}
\end{equation}
with $a = 1.28809776$ and $b = 1.31699416$.
This results in a closer match to the reference.

%

In the following subsections we will describe the statistical atomic operators \textcite{belcour2018} developed to approximate the specific light transport interactions reflection, refraction, absorption and scattering for the isotropic case.

\subsubsection{Reflection by a rough interface}
\label{sec:rough_reflection}

\textbf{Energy.} The reflection by rough surfaces is described by the rendering equation using the GGX BRDF as follows:

\begin{equation} \label{eq:rendering_equation}
    L_o(\omega_o) = \int_{\Omega} \rho_{GGX}(\omega_i, \omega_o) L_i(\omega_i) \dotproduct{\omega_i}{n} \,d\omega_i 
\end{equation}

where the subscript denotes $o$ outgoing, the subscript $i$ incoming, $L$ is the radiance, $\Omega$ is the upper hemisphere, $h$ is the half-vector. The amount of energy reflected by a surface is the directional albedo defined by:

\begin{equation}
    \FGD = \int_\Omega \rho_{GGX} \dotproduct{\omega_i}{n} \,d\omega_o
\end{equation}

Decoupling the directional albedo from the incident radiance, the energy of the outgoing radiance can be approximated:

\begin{equation}
    \int_\Omega L_o(\omega_o) \,d\omega_o \simeq \FGD \times \int_\Omega L_i(\omega_i) \,d\omega_i
\end{equation}

FGD can be precomputed parametrized by the elevation, roughness and complex index of refraction.

%

As the standard microfacet models like GGX do not account for multiple scattering energy conservation is not given. However, using a scaling factor incorporated in the FGD precomputation can compensate well for the lost energy \parencite[see e.g.][]{heitz2016,turquin2019}.

%

If the precomputation of the FGD texture is undesirable or saving the texture is too expensive, FGD can be approximated by only using the Fresnel term.\\

%

\textbf{Mean.} For high roughness values the mean gets shifted towards the normal when reflected by a rough surface. For approximating the BSDF lobe \textcite{belcour2018} ignored this shift and the outgoing mean is just the reflected incoming mean.\\

\textbf{Variance.} Reflection increases the variance. As mentioned in Section \ref{sec:interface_statistics} this increase is not linear with respect to the GGX roughness, therefore the change of  variance is applied in the introduced linear variance space.\\

\textbf{Summary.} The reflection by a rough surface of an incident energy $e_i$, mean $\mu_i$ and variance $\sigma_i$ is approximated by:
\begin{equation}
    \begin{aligned}
        e_r &= e_i \times \FGD\\
        \mu_r &= \reflect(\mu_i)\\
        \sigma_r &= \sigma_i + f(\alpha)
    \end{aligned}
\end{equation}

\subsubsection{Refraction by a rough interface}
\label{sec:rough_refraction}

\textbf{Energy.} As a consequence of conservation of energy the transmitted energy has to be $1 - \FGD$.\\

\textbf{Mean.} Snell's law defines the refraction of the incident direction. However, \textcite{belcour2018} found that the mean of the refracted lobe is different from the refracted incident direction. For high roughnesses, at grazing angles and higher refractive indices ratios deviation increases. Nevertheless, the approximation only using the refracted incident direction is good enough.\\

\textbf{Variance.} The refraction by a rough interface increases the variance depending on the roughness and bends light rays following Snell's law, leading to a scaling of the incident variance by the index of refraction ratio. The transmitted lobe can be approximated by a reflection from underneath the surface using a scaled roughness (see Figure \ref{fig:faketransmission}). 
\textcite{belcour2018} derived the following scaling factor:

\begin{equation}
    s = \frac{1}{2} \left( \frac{\eta_{12} \dotproduct{\omega_i}{n} + \dotproduct{\omega_t}{n}}{\dotproduct{\omega_t}{n}} \right)
\end{equation}

\begin{figure}[!htb]
    \includegraphics[width = \columnwidth, keepaspectratio]{fromBelcour/fake_transmission.png}
    \caption{Transmission can be approximated using a scaled reflection from underneath the interface. Image source: \textcite{belcour2018}}
    \label{fig:faketransmission}
\end{figure}

\textbf{Summary.} The refraction by a rough surface of an incident energy $e_i$, mean $\mu_i$ and variance $\sigma_i$ is approximated by:
\begin{equation} \label{eq:rough_refraction}
    \begin{aligned}
        e_t &= e_i \times (1 - \FGD)\\
        \mu_t &= \refract(\mu_i)\\
        \sigma_t &= \frac{\sigma_i}{\eta_{12}} + f(s \times \alpha)
    \end{aligned}
\end{equation}

\subsubsection{Absorption and Scattering in participating media}
\label{sec:absorption_scattering}

The effects of absorption and scattering are controlled by the depth $h$ of the layer as well as the absorption cross-section $\sigma_a$ and scattering cross-section $\sigma_s$, forming the transmittance cross-section $\sigma_t = \sigma_a + \sigma_s$.\\

\noindent\textbf{Absorption.}

\textit{Energy.} The radiance attenuation by a homogenous medium according to Beer-Lambert's law is:

\begin{equation}
    L_o(\omega_o) = \exp{\left(-\sigma_t \frac{h}{\dotproduct{\omega_o}{n}}\right)} L_i(\omega_o)
\end{equation}
where $\frac{h}{\dotproduct{\omega_o}{n}}$ is the optical depth of the layer. The energy attenuation can be approximated by evaluating the attenuation only in the incident mean direction.\\

\textit{Mean and Variance.} Even though the mean and variance are affected by absorption, \textcite{belcour2018} empirically found out that for the approximation the impact can be ignored.\\

\noindent\textbf{Scattering.} There are many aspects of scattering possible to consider. Out-scattering reduces the radiance as light collides with particles inside the medium scattering it in different directions. In-scattering on the other hand increases the radiance as light from other directions entering the medium gets scattered in the evaluated direction. The phase function describes the angular distribution of the scattered radiance. The Henyey-Greenstein phase function has an anisotropy parameter $g$ to model the amount of back-scattering, where light is mostly scattered back towards the incident direction, and forward-scattering. The approach of \textcite{belcour2018} only considers optically thin homogenous media without emission, back-scattering and multiple scattering.\\

\textit{Energy.} The amount of radiance scattered by a medium, accounting for in- and out-scattering, is given by \parencite{belcour2018,pharr2010}:
\begin{equation}
    \begin{aligned}
        L_o(\omega) &= \sigma_s \int_{0}^{h} \exp{\left(-\frac{\sigma_t t}{\dotproduct{\omega}{n}}\right)} L_s(t,\omega) \,dt \\
        \text{with}\\
        L_s(t,\omega) &= \int_{\mathcal{S}} p(\omega_i, \omega) \exp{\left(-\frac{\sigma_t (h-t)}{\dotproduct{\omega_i}{n}}\right)} L_i(\omega_i) \,d\omega_i
    \end{aligned}
\end{equation}
where $L_s(t,\omega)$ is the source function defining how much radiance is added per differential distance, accounting all incident light over the sphere $\mathcal{S}$, and $p(\omega_i, \omega)$ is the phase function.
As for absorption only the mean incident direction will be evaluated resulting in the following approximation for the outgoing energy:
\begin{equation}
    \int_{\Omega} L_o(\omega) \,d\omega \approx \sigma_s h \exp{\left(-\frac{\sigma_t t}{\dotproduct{\omega}{n}}\right)} \int_{\Omega} L_i(\omega) \,d\omega
\end{equation}

\textit{Mean.} For forward-scattering, which is only accounted for in the approximation, the impact on the mean is negligible.\\

\textit{Variance.} The width of the phase function increases the incident variance. Using the Henyey-Greenstein phase function \textcite{belcour2018} found a mapping from its anisotropy factor $g$ to GGX roughness in the linear variance space:
\begin{equation}
    \sigma_g = \left(\frac{1 - g}{g}\right)^{0.8} \frac{1}{1+g}
\end{equation}

\textbf{Summary.} The transmission through participating media of an incident energy $e_i$, mean $\mu_i$ and variance $\sigma_i$ is approximated by:
\begin{equation}
    \begin{aligned}
        e_t &= e_i \times \left[ \left(1 +  \frac{\sigma_s h}{\dotproduct{\omega_o}{n}}\right) \times \exp{\left(-\frac{\sigma_t h}{\dotproduct{\omega_o}{n}}\right)} \right]\\
        \mu_t &= \mu_i\\
        \sigma_t &= \sigma_i + \sigma_g \times \frac{\sigma_s}{\sigma_t}
    \end{aligned}
\end{equation}

\subsection{Statistics of multiple layers}
\label{sec:layered_statistics}

In this section we will present the new adding-doubling method \textcite{belcour2018} introduced. It is used to evaluate multiple scattering in plane parallel media aggregating the infinite possible light paths within it. Applying it iteratively for each layer on the incident energy and variance, the directional statistics of the individual layers get combined into an estimation of the reflection and transmission statistics of the stack of layers (see Figure \ref{fig:addingdoubling} for an illustration). As shown in Section \ref{sec:interface_statistics} the mean is approximated only using the reflected or refracted direction so the adding-doubling is not applied on it.

\begin{figure*}[h]
    \includegraphics[width = \textwidth, keepaspectratio]{fromBelcour/addingdoubling.png}
    \caption{An illustration how the presented method approximates light transport in layered materials. All possible light paths through a specific stack of layers is approximated by its statistics (leftmost). The atomic operators defined in Section \ref{sec:interface_statistics} are used to model light interactions with an interface (left). The adding-doubling method combines the statistics of all possible light paths between two interfaces (right). After iterating over all layers, the evaluated statistics can be used to instantiate BSDF lobes (rightmost). Image source: \textcite{belcour2018}.}
    \label{fig:addingdoubling}
\end{figure*}

%

\textbf{Adding-Doubling for Energy.} All operators defined in Section \ref{sec:interface_statistics} are affine transformations, therefore multiple scattering between two interfaces can be expressed using a discrete form. For that the reflection and transmission coefficients of radiance with respect to the propagation direction are needed. For example $r_{12}$ is the reflection coefficient of the interface for radiance propagating downward from medium 1 to medium 2 (see also Figure \ref{fig:coefficients}).
\begin{figure}[!htb]
    \includegraphics[width = \columnwidth, keepaspectratio]{fromBelcour/coefficients.png}
    \caption{The reflectance $r$ and transmission $t$ coefficients of one interface (left) and a stack of two interfaces (right). $l_i^-$ denotes the incident light from above the interface and $l_i^+$ from underneath. Image source: \textcite{belcour2018}.}
    \label{fig:coefficients}
\end{figure}
Within two layers all inter-reflections ($TR^+T$ paths) need to be considered when evaluating the light transport.

%

The upward reflectance between two layers accounting for multiple scattering is:
\begin{equation}
    r_{13} = r_{12} + \sum_{k=0}^{\infty} t_{12} r_{23}^{k+1} r_{21}^{k} t_{21} 
\end{equation}
This arithmetic series has the following closed form:
\begin{equation}
    r_{13} = r_{12} + \frac{t_{12} r_{23} t_{21}}{1 - r_{23} r_{21}}
\end{equation}
Similar forms can be derived for the downward reflectance, upward transmission and downward transmission (see Appendix \ref{sec:appendix_a}). The adding algorithm iteratively combines the layer statistics evaluating the reflectance $r_{13}$, $r_{31}$ and transmittance $t_{13}$, $t_{31}$. For each layer the already combined reflectance and transmittance of all layers above the current layer will be used treating it as one virtual interface.

%

The doubling algorithm duplicates a layer generating a virtual layer twice as thick. Therefore, the adding algorithm will be evaluated with $r_{12} = r_{23}$, $t_{12} = t_{23}$, $r_{21} = r_{32}$ and $t_{21} = t_{32}$. Applied multiple times on a layer of very small depth, multiple scattering inside a homogenous medium of any depth can be accomplished.

%

The reflected and transmitted energy for each of the layers in the stack can be evaluated by the adding-doubling method using the terms described in Section \ref{sec:interface_statistics}:
\begin{itemize}
    \item{
        Rough interface:
        \begin{equation}
            \begin{aligned}
                r_{12} &= r_{21} = \FGD\\
                t_{12} &= t_{21} = 1-\FGD
            \end{aligned}
        \end{equation}
    } 
    \item{
        Participating medium:
        \begin{equation}
            \begin{aligned}
                r_{12} &= r_{21} = 0\\
                t_{12} &= t_{21} = \left[ \left(1 +  \frac{\sigma_s h}{\dotproduct{\omega_o}{n}}\right) \times \exp{\left(-\frac{\sigma_t h}{\dotproduct{\omega_o}{n}}\right)} \right]
            \end{aligned}
        \end{equation}
    }
\end{itemize}

\textbf{Adding-Doubling for Variance.} As for adding-doubling for energy, the operators on variance defined in Section \ref{sec:interface_statistics} are affine transformations making the adding-doubling methodology applicable. For the approximation the variance of the average reflectance and transmission, being the variance of a weighted sum, is used. The property, that the variance of a weighted sum of distributions is the weighted sum of the distributions individual variances, allows the average reflected variance accounting for multiple scattering to be written as:
\begin{equation}
    \begin{aligned}
        \tilde{\sigma}_{13}^R = &r_{12} \sigma_{12}^R \\
        &+ \sum_{k=0}^{\infty} \left(t_{12} r_{23}^{k+1}r_{21}^k t_{21}\right) \left[\sigma_{12}^T \left(\sigma_{23}^R + k \sigma_{23}^R \sigma_{12}^R\right) J_{21} + \sigma_{21}^T\right]
    \end{aligned}
\end{equation}
where $\sigma_{ij}^R$ is the variance that is added when reflecting on the interface between medium $i$ and medium $j$, respectively $\sigma_{ij}^T$ when transmitting, and $J_{ij}$ is the transmission scaling factor. Splitting this into a geometric series and an arithmetico-geometric series gives the following close form:
\begin{equation}
    \begin{aligned}
        \tilde{\sigma}_{13}^R = &r_{12} \sigma_{12}^R \\
        &+ \left(\frac{t_{12} r_{23} t_{21}}{1 - r_{23} r_{21}}\right) \left[\sigma_{12}^T + \sigma_{21}^T + J_{12} \left(\sigma_{23}^R + \frac{r_{23} r_{21}}{1 - r_{23} r_{21}}\sigma_{12}^R\right)\right]
    \end{aligned}
\end{equation} 
This average variance needs to be normalized for the next iteration of the adding-doubling algorithm:
\begin{equation}
    \sigma_{13}^R = \frac{\tilde{\sigma}_{13}^R}{r_{12}+\left(\frac{t_{12} r_{23} t_{21}}{1 - r_{23} r_{21}}\right)}
\end{equation}
Similar forms can be derived for the average transmitted variance (see Appendix \ref{sec:appendix_a}). $\sigma_{13}^R$, $\sigma_{31}^R$, $\sigma_{13}^T$ and $\sigma_{31}^T$ will be used in place of $\sigma_{12}^R$, $\sigma_{21}^R$, $\sigma_{12}^T$ and $\sigma_{21}^T$ for the next iteration. The transmission scaling factors $J_{12}$ and $J_{21}$ also need to represent the combined transmission of the virtual interface of the previously combined layers in the next iteration, therefore the product of them will be used. For the coefficients of the current layer the terms described in Section \ref{sec:interface_statistics} will be used:
\begin{itemize}
    \item{
        Rough interface:
        \begin{equation}
            \begin{aligned}
                \sigma_{23}^R &= \sigma_{32}^R = f(\alpha)\\
                \sigma_{23}^T &= f(s \times \alpha)\\
                \sigma_{32}^T &= f(s \times \alpha)\\
                J_{23} &= \frac{1}{\eta_{12}}\\
                J_{32} &= \frac{1}{\eta_{21}}
            \end{aligned}
        \end{equation}
    } 
    \item{
        Participating medium:
        \begin{equation}
            \begin{aligned}
                \sigma_{23}^R &= \sigma_{32}^R = 0\\
                \sigma_{23}^T &= \sigma_{32}^T = \sigma_g \times \frac{\sigma_s}{\sigma_t}\\
                J_{23} &= J_{32} = 1
            \end{aligned}
        \end{equation}
    }
\end{itemize}

\subsection{BSDF}
\label{sec:bsdf}

The resulting BSDF model is a weighted sum of GGX lobes. First, the statistics of the stack of $N$ layers need to be determined using the adding-doubling method. As initial values the energy $e_i = 1$, the mean $\mu_i = \begin{pmatrix}\omega_{i_x} & \omega_{i_y}\end{pmatrix}$ and the variance $\sigma_i = 0$ will be used. For each layer $k$ one GGX lobe will be instantiated using the just evaluated vector of statistics $\begin{pmatrix}e_k & \mu_k & \sigma_k\end{pmatrix}_{k=0,\dots,N}$:
\begin{itemize}
    \item Fresnel $F$ with energy $e_k$
    \item incident direction with $\omega_{i_k} = \reflect(\mu_k)$
    \item roughness with $\alpha_k = f^{-1}(\sigma_k)$
\end{itemize}
For the BRDF this lead to:
\begin{equation}
    \rho(\omega_i, \omega_o) = \sum_{k=0}^{N} \frac{e_k G(\omega_{i_k}, \omega_o) D(\frac{\omega_{i_k}+\omega_o}{\left\lVert\omega_{i_k}+\omega_o\right\rVert })}{4 \dotproduct{\omega_{i_k}}{n} \dotproduct{\omega_o}{n}}
\end{equation}

%

In order to get a \textit{symmetric model}, \textcite{belcour2018} suggests using the half vector $h$ instead of $n$ to evaluate the adding-doubling algorithm. For \textit{importance sampling} one lobe can be randomly selected based on the energy and the VNDF (visual normal distribution function) of that lobe can then be sampled. For better results multiple importance sampling on different lobes can be used.

\subsection{Anisotropic Interfaces}
\label{sec:anisotropic_interfaces}

The model of \textcite{belcour2018} presented in the previous sections only supports isotropic interfaces. \textcite{weier2020} and \textcite{yamaguchi2019} extended the model to support anisotropic interfaces. Both of them build upon the insight, that when an anisotropic BSDF lobe is projected to the $(x,y)$ plane in tangent space the principal directions of the distribution align with the tangent and binormal direction. This property holds even for incident directions not aligned with the axes of the tangent frame or opposing anisotropies. Therefore, the projected distribution can be approximated using an elliptical distribution. This allows to evaluate the variance for each anisotropy axis independently. Thus, the adding-doubling algorithm is evaluated using a two-component vector $\begin{pmatrix}\sigma_x & \sigma_y\end{pmatrix}$ instead of a scalar for the variance. 

%

\textcite{yamaguchi2019} further extended using a covariance matrix to represent the variance of a lobe:
\begin{equation}
    \Sigma = 
    \begin{pmatrix}
        t_x & t_y    
    \end{pmatrix}^T
    \begin{pmatrix}
        \sigma_x & 0\\
        0 & \sigma_y
    \end{pmatrix}
    \begin{pmatrix}
        t_x & t_y    
    \end{pmatrix}
\end{equation}
where $t_x$ and $t_y$ are the anisotropy axes defined by the tangent and binormal. This adds the ability to rotate the tangent vector field for each layer, which makes it possible to model even more complex materials. For instantiating the resulting lobes the eigenvalues of the covariance matrix of each layer need to be determined as they are the variance of the layer.\\

%

As the energy $e$ and mean $\mu$ are negligibly affected by the anisotropy, they can be evaluated as presented in the previous sections. However, in order to precompute FGD two more dimensions would be needed for the azimuth angle and the second roughness. \textcite{weier2020} proposed an alternative approach using the split-sum approach of \textcite{karis2013}.

%

The \textit{split-sum approximation} on the one hand pre-filters the environment map for different roughness values sampling the GGX distribution and on the other hand precomputes the directional albedo using Schlick's Fresnel. This way, the precomputed directional albedo is only parametrized by roughness and elevation. To use this 2D texture of the isotropic case for the anisotropic, the incident direction $\omega_i$ needs to be stretched by the difference in roughness when fetching the texture. The elevation $\omega_z$ used to fetch the texture should be:
\begin{equation}
    \begin{aligned}
        \omega_z &= \sqrt{1-\left(\omega_x^2 + \omega_y^2\right)}\\
        \text{with}\\
        \omega_x &= \frac{\omega_{i_x}}{1+\left(\Delta \alpha_x\right)^2}\\
        \omega_y &= \frac{\omega_{i_y}}{1+\left(\Delta \alpha_y\right)^2}\\
        \Delta \alpha_x &= \alpha_x - \min{\alpha_x, \alpha_y}\\
        \Delta \alpha_y &= \alpha_y - \min{\alpha_x, \alpha_y}\\
    \end{aligned}
\end{equation}
The roughness used to fetch the texture should be $\min{\alpha_x, \alpha_y}$. This method does not account for multiple scattering, but like for $\FGD$ mentioned earlier it can be rescaled.

%

For \textit{image based lighting} (IBL) of anisotropic BRDF's there exists no real solution for pre-filtering environment maps so far. Therefore, the usage for real-time rendering of anisotropic layered materials is limited as Monte Carlo integration is needed. Monte Carlo integration with importance sampling estimates an integral $I = \int_\Omega f(x) \,dx$ drawing $N$ samples $x_i$ from a distribution $p(x)$ as follows:
\begin{equation}
    I \approx \frac{1}{N} \sum_{i=1}^{N} \frac{f(x_i)}{p(x_i)}
\end{equation}
For IBL $N$ light samples of the environment map can be drawn to approximate $\int_{\Omega} L_i(\omega_i) \,d\omega_i$. Therefore, normals can be sampled from the VNDF of GGX \parencite{heitz2018}. Using them, the light direction can be computed to fetch the environment texture. To reduce the sample count needed, mipmaps for the environment map can be generated and depending on the PDF of the current sample the according mipmap level can be chosen \parencite{nguyen2007}. For a low PDF of a sample the probability of other generating samples in a similar direction is low. Therefore, the sample should account for a larger area of the environment map to provide a better approximation of the incident light, leading to a higher mipmap level. The mipmap level $l$ a sample should fetch is:
\begin{equation}
    l = \max{\frac{1}{2} \log_2{\left(\frac{\Omega_s}{\Omega_t}\right)},0}
\end{equation}
where $\Omega_s$ is the solid angle of the sample and $\Omega_t$ is the solid angle of a texel at mipmap level 0.